%% Research Note of the AAS: the transiting companion of Boyajian's star. 2026-11-24.
%% Text from the Claude Doc "RNAAS draft: Boyajian's star companion period" (rev 20, Notes section removed).
%% Source, numbers and provenance: https://github.com/coracle-us/duo-transits (boyajian-companion/; see rnaas/numbers.md).
\documentclass{aastex701}

\begin{document}

\title{Kepler Rules Out Most $\sim$780-Day Periods and Window 4 for the Transiting Companion of Boyajian's Star}

\author[orcid=0009-0000-7524-4470]{Jon McBee}
\affiliation{AFFILIATION PLACEHOLDER}
\email{jon.mcbee@gmail.com}

\begin{abstract}
A symmetric, 1.1\%-deep, $\sim$21\,h single transit of KIC~8462852 in TESS Sector~15 has been attributed to a
$\sim$9\,$M_{\rm J}$ companion with five Kepler-allowed period windows between 1030 and 1368\,d, while a recurrence
model of the star's dimming families favors $\sim$780\,d. I test trial periods from 770 to 845\,d, and the five
windows, against Kepler long-cadence photometry under one constraint: flux deficits add, so a transit cannot hide
in data dimmed by less than the transit itself. Near 780\,d, Kepler allows only 776.17--776.60, 790.50--790.66,
792.08--792.92 and 793.25--793.40\,d, excluding the model's maximum-likelihood (781.6\,d) and median (800.9\,d)
periods. The same test excludes window~4 (1345.65\,d), whose earlier transit falls on a dip only 0.45\% deep.
Public AAVSO and ASAS-SN photometry can neither detect nor exclude a transit in window~3 (2026 October 18--23),
the published favorite.
\end{abstract}

\keywords{\uat{Exoplanet detection methods}{489} --- \uat{Transit photometry}{1709}}

\section{Introduction}

Boyajian's star (KIC~8462852) shows deep, aperiodic, chromatic dips usually attributed to dust from disrupted
exocomets or planetesimals \citep{Boyajian2016}. \citet[hereafter MF26]{MF26} found a single symmetric transit
in TESS Sector~15 (2019 September~3; depth $\sim$1.1\%, duration $\sim$21\,h). With SOPHIE radial velocities
(RVs), it implies a $9.4^{+4.9}_{-4.1}\,M_{\rm J}$ companion with $P=1211^{+89}_{-78}$\,d. Assuming any earlier
transit fell in a Kepler gap or under a dip deeper than 0.1\%, MF26 allow five windows near 1030, 1035, 1302,
1345 and 1368\,d, and rank window~3 (1301.0--1303.6\,d) most probable. \citet{Zhu2026} fit the same transit with
a ringed Jupiter-sized planet.

\citet[hereafter SN26]{SN26} instead models the dimming families as returns of a disrupted population, with
maximum likelihood at 781.6\,d and posterior median $800.9^{+41.1}_{-28.4}$\,d. SN26 suggests the TESS object
may be the surviving remnant, noting that the D800-to-TESS interval is four orbits of 776.1\,d. I test which of
these periods Kepler still allows.

\section{Data and method}

I used Kepler Q0--Q17 long-cadence PDCSAP photometry \citep[][SAP\_QUALITY $=0$, each quarter normalized to its
median]{Borucki2010} and TESS SPOC 2-minute photometry \citep{Ricker2015} from Sectors 14, 15, 41, 54 and 55, all
from MAST. A trapezoid fit to the Sector~15 event gives mid-transit BTJD $1730.154\pm0.003$, depth 1.16\% and
total duration 20.3\,h, within $2\sigma$ of MF26's $T_0=$ BJD $2458730.164\pm0.004$, which I adopt.

For each trial period $P$ I predict $t_n=T_0+nP$ within the Kepler baseline and take the local quiet level as
the 95th percentile of flux within $\pm3$\,d. Because dust and companion dim the star together, a real transit
leaves every cadence in its core ($|t-t_n|<0.25$\,d) at least $\sim$1\% below that level. I exclude $P$ if any
predicted core holds three or more cadences dimmed by less than 0.8\%; cores in gaps, or inside dips at least
0.8\% deep, survive. Core half-widths of 0.20--0.25\,d, or my own $T_0$, shift the surviving ranges by
$\leq0.02$\,d. As a check, the code recovers MF26's windows 1, 2, 3 and 5 to within $\sim$0.5\,d at each edge.

\section{Results}

\textbf{The $\sim$780\,d hypothesis.} Between 770 and 845\,d, the $n=-4$ and $n=-3$ transits fall in Kepler
Quarters 7--9 and 15--17 (Figure~\ref{fig:alias}). Four narrow ranges survive. At 776.17--776.60\,d, both earlier
transits land inside the two deepest events, D800 and the final D1500 dip; at 790.50--790.66\,d, both fall in
gaps; at 792.08--792.92 and 793.25--793.40\,d, $n=-4$ falls in the Q7/Q8 gap and $n=-3$ in the deepest D1500
dip. SN26's 781.6 and 800.9\,d periods are excluded: their earlier transits land on quiet data with
0.02--0.07\% scatter, where a 1.1\% transit would stand out at $>30\sigma$. Even 776.1\,d is marginally excluded,
because its $n=-3$ core falls on the egress of the last D1500 dip, within 0.6--0.8\% of quiet.

Every surviving range needs a transit hidden in a gap or in one of the deepest dips. That is allowed, but not
evidence. If the TESS object is SN26's remnant, its period must be 776.2--776.6\,d, plausible only if D800 and
the final D1500 dip came from material near the companion.

\textbf{The $\gtrsim$1030\,d bound.} MF26's $\sim$1030\,d floor, repeated by \citet{Zhu2026}, is not a Kepler
limit. It comes from their TESS plus RV posterior, which rests on a $2.3\sigma$ semi-amplitude; Kepler alone
excludes only $P<776$\,d, and their TESS-only fit gives $1103^{+242}_{-234}$\,d. The 776.2--776.6\,d range is
therefore disfavored by the RVs, not excluded. MF26 report no transit on the night of 2023 December~3 (BJD
2460282.2--2460282.4), but that night covers one only for $P\leq776.31$\,d, leaving 776.3--776.6\,d untested.

\textbf{Window 4.} For $P=1345.648\pm0.008$\,d, the $n=-2$ transit falls at BKJD 1205.9, on a 0.45\%-deep dip,
not a gap. MF26's 0.1\% rule masks it, but a 1.1\% transit cannot produce a shallower dip. Window~4 is excluded
unless the transit is much shallower than in TESS. Windows 1, 2, 3 and 5 survive.

\section{Window 3 (2026 October)}

MF26's most probable window predicted a transit between 2026 October 18 and 23. No TESS sector covers it; the
next visits are expected in Sectors 120--121 (2027 August--September). I tested it with AAVSO photometry
\citep[22,234 observations from 2021 July to 2026 October 29 after cuts;][]{AAVSO} and ASAS-SN
\citep{Kochanek2017,Hart2023}, comparing each observer and band with its own median over $\pm15$\,d. A trial
mid-time is excluded if a 1.1\% transit raises $\chi^2$ by more than 25 over the core points, and supported if it
lowers $\chi^2$ by more than 25.

Window~3 remains untested. AAVSO holds 19 single visits in the window, on six nights from three observers. Only
43\% of trial mid-times put any point in the core, and none reaches either threshold ($\Delta\chi^2$ from $-3.4$
to $+13.3$); night-to-night scatter of 7--24\,mmag is comparable to the $\sim$12\,mmag depth. The ASAS-SN light
curve ends in 2025 June. Earlier returns of the surviving $\sim$780\,d ranges (2021 October--November, 2023
December, 2024 January) and of windows 3 and 5 (2023 March and June) are also unconstrained; the closest, window
3's 2023 March return, reaches $\Delta\chi^2=+19.4$ on 15\% of its mid-times. The 2026 returns of the $\sim$780\,d
ranges fell at solar conjunction or in morning twilight.

Gaia DR4 epoch astrometry (2026 December), suggested by both MF26 and \citet{Zhu2026}, can constrain the orbit
without another transit.

\begin{figure}[ht!]
\plotone{boyajian_alias_figure.png}
\caption{Top: Kepler status of the predicted $n=-4$ and $n=-3$ transits for trial periods of 770--812\,d: blue,
quiet data that exclude a 1.1\% transit; orange, a dip at least 0.8\% deep; grey, a gap. Red marks surviving
periods. The shaded band is SN26's 68\% interval, the dashed line their maximum-likelihood period and the dotted
line 776.1\,d. Bottom: Kepler flux at predicted epochs. The 781.63\,d and window~4 epochs (top row) are compared
with the expected transit (red); for 776.30\,d (bottom row) the predicted transit (shaded) falls inside D800 and
the final D1500 dip.\label{fig:alias}}
\end{figure}

\begin{acknowledgments}
I acknowledge with thanks the variable star observations from the AAVSO International Database contributed by
observers worldwide and used in this research. This work uses Kepler
\dataset[doi:10.17909/T9488N]{https://doi.org/10.17909/T9488N} and TESS SPOC
\dataset[doi:10.17909/t9-nmc8-f686]{https://doi.org/10.17909/t9-nmc8-f686} light curves from MAST.
\end{acknowledgments}

\facilities{Kepler, TESS, AAVSO, ASAS-SN}

\section*{AI disclosure statement}
I used Anthropic's Claude, an AI assistant, under my supervision to write and run the analysis code and to draft
and edit this note. I reviewed the text, checked every number against the saved outputs and every reference
against ADS, and take full responsibility for the content. The figure was made by the analysis code from the
data; the note contains no AI-generated figures.

\begin{thebibliography}{}
\bibitem[AAVSO(2026)]{AAVSO} American Association of Variable Star Observers (AAVSO) 2026, AAVSO International Database
\bibitem[Borucki et al.(2010)]{Borucki2010} Borucki, W.~J., Koch, D., Basri, G., et al. 2010, Science, 327, 977
\bibitem[Boyajian et al.(2016)]{Boyajian2016} Boyajian, T.~S., LaCourse, D.~M., Rappaport, S.~A., et al. 2016, MNRAS, 457, 3988
\bibitem[Hart et al.(2023)]{Hart2023} Hart, K., Shappee, B.~J., Hey, D., et al. 2023, arXiv:2304.03791
\bibitem[Kochanek et al.(2017)]{Kochanek2017} Kochanek, C.~S., Shappee, B.~J., Stanek, K.~Z., et al. 2017, PASP, 129, 104502
\bibitem[Madurga-Favieres et al.(2026)]{MF26} Madurga-Favieres, C., Str{\o}m, P.~A., Wilson, T.~G., Heidari, N., \& Kiefer, F. 2026, MNRAS, 550, stag1273
\bibitem[Ricker et al.(2015)]{Ricker2015} Ricker, G.~R., Winn, J.~N., Vanderspek, R., et al. 2015, JATIS, 1, 014003
\bibitem[Socas-Navarro(2026)]{SN26} Socas-Navarro, H. 2026, arXiv:2609.02726
\bibitem[Zhu \& Liu(2026)]{Zhu2026} Zhu, W., \& Liu, Q. 2026, arXiv:2610.34806
\end{thebibliography}

\end{document}
